% !TeX root = part9.tex
\documentclass[../final.tex]{subfiles}
\begin{document}
\ifSubfilesClassLoaded{\setcounter{chapter}{6}}{}

% Sources: Desktop/1.jpg, 2.jpg, 3.jpg.
\chapter{Тепловое расширение и эффект Мёссбауэра}
\rsubtitle{Закон Грюнейзена, отдача ядра и бесфононные переходы}

\section{Соответственные состояния при заданном давлении}
Модель Дебая связывает температурные зависимости колебательных величин
с одним энергетическим масштабом $\Theta_D=\hbar\omega_D$.
Температуру далее измеряем в энергетических единицах: $T=k_BT_{\mathrm K}$.
Число атомов $N$ фиксировано.

Для потенциала Гиббса используем приближение одного частотного масштаба:
\begin{equation*}
 \Phi(p,T)\simeq\Phi_0(p)+N\Theta(p)f(x),\qquad
 x=\frac{T}{\Theta(p)},\qquad \Theta\equiv\Theta_D.
\end{equation*}
Это запись закона соответственных состояний при заданном давлении.
Она применима, когда тепловое изменение объёма мало и изменение спектра
можно описать зависимостью $\Theta(p)$. Функция $f$ безразмерна,
а $\Phi_0$ содержит температурно независимую часть, включая нулевые колебания.

\subsection{Энтропия и теплоёмкость}
Из дифференциала $d\Phi=-S\,dT+V\,dp$ следует
\begin{align*}
 S&=-\left(\frac{\partial\Phi}{\partial T}\right)_p
 =-N\Theta f'(x)\frac1\Theta=-Nf'(x),\\
 C_p&=T\left(\frac{\partial S}{\partial T}\right)_p
 =-\frac{NT}{\Theta}f''(x)=-Nx f''(x).
\end{align*}
Штрихи у $f$ означают производные по $x$.
При $C_p>0$ имеем $f''(x)<0$.

\subsection{Объём как производная потенциала}
При фиксированной температуре меняются и множитель $\Theta$, и аргумент $x$:
\begin{equation*}
 \Theta'=\frac{d\Theta}{dp},\qquad
 \left(\frac{\partial x}{\partial p}\right)_T
 =-\frac{T\Theta'}{\Theta^2}=-\frac{x\Theta'}\Theta.
\end{equation*}
Поэтому
\begin{align*}
 V&=\left(\frac{\partial\Phi}{\partial p}\right)_T
 =\Phi_0'(p)+N\Theta'f(x)
 +N\Theta f'(x)\left(-\frac{x\Theta'}\Theta\right)\\
 &=V_0(p)+N\Theta'\bigl[f(x)-xf'(x)\bigr],
 \qquad V_0(p)=\Phi_0'(p).
\end{align*}

\clearpage
\section{Тепловое расширение и закон Грюнейзена}
Объёмный коэффициент теплового расширения определён как
\begin{equation*}
 \alpha_V=\frac1V\left(\frac{\partial V}{\partial T}\right)_p.
\end{equation*}
Дифференцируем найденный объём, сохраняя все члены:
\begin{align*}
 \left(\frac{\partial V}{\partial T}\right)_p
 &=N\Theta'\left[\frac{f'(x)}\Theta
 -\frac{f'(x)}\Theta-\frac{x f''(x)}\Theta\right]\\
 &=-\frac{N\Theta'}\Theta x f''(x)
 =\frac{\Theta'}\Theta C_p.
\end{align*}
Члены с $f'$ сокращаются. В принятом приближении получаем закон Грюнейзена:
\begin{equation*}
 \boxed{\quad \frac{\alpha_V}{C_p}
 =\frac1{V\Theta}\frac{d\Theta}{dp}
 \simeq\frac1{V_0\Theta}\frac{d\Theta}{dp}.\quad}
\end{equation*}
Последний переход использует малость теплового изменения объёма.
Если $d\Theta/dp>0$, нагревание приводит к расширению.
При слабой температурной зависимости правой части $\alpha_V(T)$
повторяет температурный ход $C_p(T)$.
Для изотропного тела линейный коэффициент равен $\alpha_L=\alpha_V/3$.
При переходе к кельвинам и $\alpha_V$, и $C_p$ умножаются на $k_B$;
их отношение сохраняется.

\subsection{Межатомный потенциал и характерная частота}
Если $a$ --- минимум парного потенциала $U(r)$, то вблизи него
\begin{equation*}
 U(r)=U(a)+\frac{U''(a)}2(r-a)^2+\cdots,\qquad
 \omega_0^2=\frac{U''(a)}m.
\end{equation*}
Здесь $m$ --- масса, соответствующая выбранной координате колебаний.
Частота Дебая задаётся всем спектром решётки и имеет тот же характерный
масштаб жёсткости связей. Увеличение жёсткости при сжатии обычно повышает
$\omega_D$ и $\Theta_D$.

\begin{rbox}[unbreakable,left=14pt,right=14pt,top=8pt,bottom=6pt]{[]}
\centering
\captionsetup{font=footnotesize,hypcap=false,skip=5pt}
\begin{tikzpicture}[>=Stealth,font=\small,x=1cm,y=0.7cm]
 \draw[->] (0,0)--(10.4,0) node[right] {$r$};
 \draw[->] (0,-2.1)--(0,1.6) node[above] {$U(r)$};
 \draw[p3,thick] (0.7,1.3) .. controls (0.85,-0.7) and (1.2,-1.8) .. (2,-1.8)
 .. controls (3,-1.8) and (3.3,-0.2) .. (9.8,-0.08);
 \draw[dashed] (2,0)--(2,-1.8);
 \node[above] at (2,0) {$a$};
\end{tikzpicture}
\captionof{figure}{Асимметричная яма: резкое отталкивание и более пологая область притяжения.}
\end{rbox}
В чисто гармонической яме среднее положение не сдвигается при нагревании.
Тепловое расширение связано с ангармонизмом; в квазигармоническом описании
он учитывается через зависимость частот от объёма или давления.

\clearpage
\section{Отдача свободного ядра}
Пусть энергия ядерного перехода равна $\varepsilon_0$.
Фотон с волновым вектором $\mathbf k$ имеет энергию и импульс
\begin{equation*}
 \varepsilon_\gamma=\hbar\omega=\hbar ck,\qquad
 \mathbf p_\gamma=\hbar\mathbf k.
\end{equation*}
Если покоившееся ядро массы $M$ испускает фотон, то по закону сохранения
импульса оно получает импульс $-\hbar\mathbf k$.
Его кинетическая энергия отдачи равна
\begin{equation*}
 \boxed{\quad R=\frac{p_\gamma^2}{2M}
 =\frac{\hbar^2k^2}{2M}=\frac{\varepsilon_\gamma^2}{2Mc^2}.\quad}
\end{equation*}
Используем нерелятивистское движение ядра, $\varepsilon_0\ll Mc^2$.
В выражении для малой поправки $R$ можно положить
$\varepsilon_\gamma\simeq\varepsilon_0$.

\subsection{Испускание и поглощение покоящимися ядрами}
При испускании часть энергии перехода уходит на отдачу:
\begin{equation*}
 \varepsilon_0=\varepsilon_\gamma^{\text{исп}}+R,
 \qquad \varepsilon_\gamma^{\text{исп}}\simeq\varepsilon_0-R.
\end{equation*}
При поглощении фотон должен обеспечить и возбуждение, и отдачу:
\begin{equation*}
 \varepsilon_\gamma^{\text{погл}}=\varepsilon_0+R.
\end{equation*}
Центры линий испускания и поглощения разнесены на $2R$.
Поэтому одинаковые свободные покоящиеся ядра не обеспечивают резонанс,
если $R$ значительно превосходит естественную ширину линии.

\subsection{Оценка для $^{57}\mathrm{Fe}$}
Для перехода с энергией $\varepsilon_0\simeq14.4\,\text{кэВ}$
и $Mc^2\simeq57\cdot931.5\,\text{МэВ}\simeq53.1\,\text{ГэВ}$
\begin{equation*}
 R\simeq\frac{(1.44\cdot10^4)^2}{2\cdot5.31\cdot10^{10}}\,\text{эВ}
 \simeq1.95\cdot10^{-3}\,\text{эВ}.
\end{equation*}
Естественная ширина связана со средним временем жизни возбуждённого уровня:
\begin{equation*}
 \Gamma=\frac\hbar\tau,\qquad
 \tau\simeq141\,\text{нс},\qquad
 \Gamma\simeq4.7\cdot10^{-9}\,\text{эВ}.
\end{equation*}
Следовательно, $R/\Gamma\simeq4.2\cdot10^5$:
энергия отдачи мала по сравнению с энергией фотона, но огромна по сравнению
с шириной ядерной линии.

\clearpage
\section{Движущееся ядро: доплеровское уширение}
До испускания ядро имеет импульс $\mathbf P=\hbar\mathbf K=M\mathbf v$;
после испускания --- $\mathbf P-\hbar\mathbf k$.
Энергии начального и конечного состояний равны
\begin{equation*}
 E_i=\varepsilon_0+\frac{P^2}{2M},\qquad
 E_f=\frac{|\mathbf P-\hbar\mathbf k|^2}{2M}.
\end{equation*}
Раскрывая квадрат, получаем
\begin{align*}
 \varepsilon_\gamma^{\text{исп}}&=E_i-E_f
 =\varepsilon_0+\frac{P^2-(P^2-2\hbar\mathbf k\cdot\mathbf P+\hbar^2k^2)}{2M}\\
 &=\varepsilon_0+\frac{\hbar\mathbf k\cdot\mathbf P}{M}-R
 =\varepsilon_0+\hbar\mathbf k\cdot\mathbf v-R.
\end{align*}
При поглощении импульс меняется с $\mathbf P$ на $\mathbf P+\hbar\mathbf k$:
\begin{equation*}
 \varepsilon_\gamma^{\text{погл}}
 =\varepsilon_0+\frac{|\mathbf P+\hbar\mathbf k|^2-P^2}{2M}
 =\varepsilon_0+\hbar\mathbf k\cdot\mathbf v+R.
\end{equation*}

\subsection{Средняя энергия и дисперсия}
В равновесном газе $\langle\mathbf v\rangle=0$.
Для компоненты $v_\parallel$ вдоль $\mathbf k$ распределение Максвелла даёт
\begin{equation*}
 w(v_\parallel)=\sqrt{\frac{M}{2\pi T}}
 e^{-Mv_\parallel^2/(2T)},\qquad
 \langle v_\parallel^2\rangle=\frac TM.
\end{equation*}
Поэтому
\begin{align*}
 \langle\varepsilon_\gamma^{\text{исп}}\rangle&=\varepsilon_0-R,
 &\langle\varepsilon_\gamma^{\text{погл}}\rangle&=\varepsilon_0+R,\\
 \sigma_\varepsilon^2
 &=\langle\varepsilon_\gamma^2\rangle-\langle\varepsilon_\gamma\rangle^2
 =\hbar^2k^2\langle v_\parallel^2\rangle
 =2RT,
 &\sigma_\varepsilon&=\sqrt{2RT}.
\end{align*}
Дисперсия растёт как $T$, а ширина --- как $\sqrt T$.
При $T_{\mathrm K}=300\,\text{К}$ для $^{57}\mathrm{Fe}$
$\sigma_\varepsilon\simeq1.0\cdot10^{-2}\,\text{эВ}$.

\begin{rbox}[unbreakable,left=14pt,right=14pt,top=7pt,bottom=6pt]{[]}
\centering
\captionsetup{font=footnotesize,hypcap=false,skip=5pt}
\begin{tikzpicture}[>=Stealth,font=\footnotesize,x=0.95cm,y=0.65cm]
 \foreach \yy/\ww/\hh in {0/0.10/1.45,2.6/1.40/1.05} {
  \draw[->] (0,\yy)--(11,\yy) node[right] {$\varepsilon_\gamma$};
  \draw[p3,thick,domain=0.4:10.4,samples=120]
   plot (\x,{\yy+\hh*exp(-((\x-3.6)^2)/\ww)});
  \draw[red3,thick,domain=0.4:10.4,samples=120]
   plot (\x,{\yy+\hh*exp(-((\x-6.8)^2)/\ww)});
  \node[below] at (3.6,\yy) {$\varepsilon_0-R$};
  \node[below] at (6.8,\yy) {$\varepsilon_0+R$};
  \draw (5.2,\yy)--(5.2,{\yy-0.08});
  \node[below] at (5.2,\yy) {$\varepsilon_0$};
 }
 \node[p3] at (2.4,4.35) {испускание};
 \node[red3] at (7.9,4.35) {поглощение};
 \node at (9.6,3.9) {$T_2>T_1$};
 \node at (9.6,1.1) {$T_1$};
\end{tikzpicture}
\captionof{figure}{Доплеровское уширение может увеличить перекрытие разнесённых линий в газе.}
\end{rbox}

\clearpage
\section{Связанное ядро и бесфононный переход}
В кристалле ядро связано с решёткой. Передача импульса может сопровождаться
изменением колебательного состояния, а может происходить без рождения
и поглощения фононов. Во втором случае возникает узкая резонансная линия:
это эффект Мёссбауэра.

\subsection{Оператор передачи импульса}
Сначала рассмотрим одномерный осциллятор с координатой $x$.
Разложим начальное состояние по собственным состояниям импульса:
\begin{equation*}
 |n_i\rangle=\int dP\,|P\rangle\langle P|n_i\rangle.
\end{equation*}
При испускании фотона каждый импульс ядра уменьшается на $\hbar k$.
Поскольку $\langle x|P\rangle\propto e^{iPx/\hbar}$, имеем
\begin{equation*}
 e^{-ikx}|P\rangle=|P-\hbar k\rangle,
 \qquad
 |\psi_{\text{после}}\rangle
 =\int dP\,|P-\hbar k\rangle\langle P|n_i\rangle
 =e^{-ikx}|n_i\rangle.
\end{equation*}
Это состояние, вообще говоря, является суперпозицией уровней осциллятора:
\begin{equation*}
 |\psi_{\text{после}}\rangle=\sum_{n_f}|n_f\rangle
 \underbrace{\langle n_f|e^{-ikx}|n_i\rangle}_{A_{n_fn_i}},
 \qquad P_{n_fn_i}=|A_{n_fn_i}|^2.
\end{equation*}

\subsection{Энергия фотона и колебательные переходы}
Предположим, что частота и положение равновесия осциллятора не меняются
при ядерном переходе. Тогда
\begin{align*}
 E_i&=\varepsilon_0+\hbar\omega\left(n_i+\frac12\right),
 &E_f&=\hbar\omega\left(n_f+\frac12\right),\\
 \varepsilon_\gamma&=\varepsilon_0+\hbar\omega(n_i-n_f).
\end{align*}
Если $n_f>n_i$, часть энергии уходит в колебания и
$\varepsilon_\gamma<\varepsilon_0$; если $n_f<n_i$, фотон получает
добавочную энергию. При $n_f=n_i$ колебательная энергия сохраняется:
\begin{equation*}
 \boxed{\varepsilon_\gamma=\varepsilon_0.}
\end{equation*}
Для решётки с набором мод вместо одного слагаемого возникает сумма:
\begin{equation*}
 \varepsilon_\gamma=\varepsilon_0+
 \sum_\nu\hbar\omega_\nu(n_{i\nu}-n_{f\nu}).
\end{equation*}
В бесфононном канале набор чисел заполнения сохраняется.
Импульс передаётся кристаллу как целому; энергия его отдачи
$\hbar^2k^2/(2M_{\text{кр}})$ пренебрежимо мала из-за большой массы
$M_{\text{кр}}$. Закон сохранения импульса при этом выполняется.

\clearpage
\section{Вероятность бесфононного перехода}
\subsection{Осциллятор в основном состоянии}
Обозначим $s^2=\langle0|x^2|0\rangle=\hbar/(2M\omega)$.
Плотность вероятности координаты в основном состоянии гауссова:
\begin{equation*}
 |\psi_0(x)|^2=\frac1{\sqrt{2\pi s^2}}
 \exp\!\left(-\frac{x^2}{2s^2}\right).
\end{equation*}
Амплитуда перехода без изменения колебательного состояния равна
\begin{align*}
 A_{00}&=\langle0|e^{-ikx}|0\rangle
 =\frac1{\sqrt{2\pi s^2}}\int_{-\infty}^{\infty}
 \exp\!\left(-\frac{x^2}{2s^2}-ikx\right)dx\\
 &=\exp\!\left(-\frac{k^2s^2}{2}\right).
\end{align*}
Тот же результат виден из разложения экспоненты:
\begin{equation*}
 A_{00}=1-ik\langle x\rangle_0-\frac{k^2}2\langle x^2\rangle_0+\cdots
 =1-\frac{k^2s^2}{2}+\cdots.
\end{equation*}
Нечётные моменты равны нулю, а сумма всех чётных моментов даёт гауссову
экспоненту. При переходе от амплитуды к вероятности показатель удваивается:
\begin{equation*}
 P_{00}=|A_{00}|^2=e^{-k^2\langle x^2\rangle_0}.
\end{equation*}

\subsection{Тепловое среднее и три измерения}
Для гармонической решётки тепловые флуктуации смещений также гауссовы.
Пусть $\mathbf u$ --- смещение излучающего ядра от положения равновесия,
а $u_\parallel$ --- его проекция на направление фотона. Тогда
\begin{equation*}
 \mathcal A_T=\left\langle e^{-i\mathbf k\cdot\mathbf u}\right\rangle_T
 =\exp\!\left[-\frac12\langle(\mathbf k\cdot\mathbf u)^2\rangle_T\right].
\end{equation*}
Вес когерентной упругой линии в макроскопической гармонической решётке
определяется фактором Лэмба--Мёссбауэра:
\begin{equation*}
 f_{\mathrm{LM}}=|\mathcal A_T|^2
 =\exp[-k^2\langle u_\parallel^2\rangle_T].
\end{equation*}
При изотропных колебаниях
\begin{equation*}
 \langle u_\parallel^2\rangle_T=\frac13\langle|\mathbf u|^2\rangle_T,
 \qquad
 \boxed{f_{\mathrm{LM}}=\exp\!\left[-\frac{k^2}{3}
 \langle|\mathbf u|^2\rangle_T\right].}
\end{equation*}
Чем меньше амплитуда колебаний, тем больше доля бесфононной линии.
Даже при $T=0$ нулевые колебания дают $f_{\mathrm{LM}}<1$.

\clearpage
\section{Средний квадрат смещения через нормальные моды}
Для одноатомной решётки массы атомов равны $M$.
Собственные состояния гармонического гамильтониана имеют вид
$|\Psi\rangle=|n_1,n_2,\ldots,n_{3N}\rangle$.
Ортогональное преобразование $\boldsymbol\chi=Z\mathbf Q$
приводит квадратичную форму к диагональному виду:
\begin{equation*}
 U_2=\frac M2\boldsymbol\chi^{\mathsf T}B\boldsymbol\chi
 =\frac M2\mathbf Q^{\mathsf T}\operatorname{diag}(\omega_\nu^2)\mathbf Q
 =\frac M2\sum_\nu\omega_\nu^2Q_\nu^2.
\end{equation*}
Поскольку $Z^{\mathsf T}Z=1$, сохраняется сумма квадратов координат:
\begin{equation*}
 \sum_{i=1}^N|\mathbf u_i|^2=\sum_\nu Q_\nu^2,
 \qquad
 \langle u^2\rangle\equiv\frac1N\sum_i\langle|\mathbf u_i|^2\rangle
 =\frac1N\sum_\nu\langle Q_\nu^2\rangle.
\end{equation*}
Для эквивалентных узлов это же среднее характеризует излучающий атом.
В изотропном случае разложение упругой амплитуды начинается с
\begin{equation*}
 \langle e^{-i\mathbf k\cdot\mathbf u}\rangle
 =1-\frac{k^2}{6}\langle u^2\rangle+\cdots
 =1-\frac{k^2}{6N}\sum_\nu\langle Q_\nu^2\rangle+\cdots.
\end{equation*}

\subsection{Почему появляется произведение по модам}
Выразим проекцию смещения излучающего атома через нормальные координаты:
\begin{equation*}
 u_\parallel=\sum_\nu c_\nu Q_\nu,\qquad
 e^{-i\mathbf k\cdot\mathbf u}
 =e^{-ik\sum_\nu c_\nu Q_\nu}
 =\prod_\nu e^{-ikc_\nu Q_\nu}.
\end{equation*}
Коэффициенты $c_\nu$ задаются собственными векторами колебаний.
Координаты разных мод коммутируют, а состояние $|\Psi\rangle$
является произведением состояний отдельных осцилляторов. Поэтому
\begin{align*}
 \mathcal A_\Psi
 &=\langle\Psi|e^{-i\mathbf k\cdot\mathbf u}|\Psi\rangle
 =\prod_\nu\langle n_\nu|e^{-ikc_\nu Q_\nu}|n_\nu\rangle\\
 &=\prod_\nu\left[1-\frac{k^2c_\nu^2}{2}
 \langle n_\nu|Q_\nu^2|n_\nu\rangle+\cdots\right].
\end{align*}
При раскрытии произведения в первом порядке по квадратам смещений
получаем сумму вкладов мод. Усреднение по эквивалентным узлам
и направлениям даёт $\overline{c_\nu^2}=1/(3N)$ и возвращает
поправку $-k^2\sum_\nu\langle Q_\nu^2\rangle/(6N)$.

\clearpage
\subsection{Матричный элемент $Q_\nu^2$}
Используем операторы рождения и уничтожения:
\begin{align*}
 Q_\nu&=\sqrt{\frac\hbar{2M\omega_\nu}}(a_\nu+a_\nu^\dagger),\\
 \langle n|Q_\nu^2|n\rangle
 &=\frac\hbar{2M\omega_\nu}
 \langle n|a_\nu^2+(a_\nu^\dagger)^2
 +a_\nu a_\nu^\dagger+a_\nu^\dagger a_\nu|n\rangle\\
 &=\frac\hbar{2M\omega_\nu}\bigl[0+0+(n+1)+n\bigr]
 =\frac\hbar{2M\omega_\nu}(2n+1).
\end{align*}
Два первых слагаемых меняют число квантов на два и потому не дают
диагонального матричного элемента.

\subsection{Усреднение по распределению Бозе}
Статистическая сумма одной моды равна
\begin{equation*}
 Z_\nu=\sum_{n=0}^\infty e^{-\hbar\omega_\nu(n+1/2)/T}
 =\frac{e^{-\hbar\omega_\nu/(2T)}}{1-e^{-\hbar\omega_\nu/T}}.
\end{equation*}
Для $q=e^{-\hbar\omega_\nu/T}$ вероятность уровня равна
\begin{equation*}
 w_n=\frac{e^{-\hbar\omega_\nu(n+1/2)/T}}{Z_\nu}
 =(1-q)q^n.
\end{equation*}
Из сумм $\sum q^n=(1-q)^{-1}$ и $\sum nq^n=q/(1-q)^2$ получаем
\begin{align*}
 \langle Q_\nu^2\rangle_T
 &=\frac\hbar{2M\omega_\nu}(1-q)\sum_{n=0}^\infty(2n+1)q^n\\
 &=\frac\hbar{2M\omega_\nu}\frac{1+q}{1-q}
 =\frac\hbar{M\omega_\nu}
 \left(\frac12+\frac1{e^{\hbar\omega_\nu/T}-1}\right).
\end{align*}
Это же выражение можно записать как
\begin{equation*}
 \langle Q_\nu^2\rangle_T=\frac\hbar{2M\omega_\nu}
 \coth\!\left(\frac{\hbar\omega_\nu}{2T}\right).
\end{equation*}
При $T\to0$ остаётся $\hbar/(2M\omega_\nu)$, а при
$T\gg\hbar\omega_\nu$ получаем классический результат
$\langle Q_\nu^2\rangle_T\simeq T/(M\omega_\nu^2)$.
В тепловом состоянии независимость мод и гауссово распределение координат
позволяют просуммировать все порядки:
\begin{equation*}
 \mathcal A_T=\prod_\nu
 \exp\!\left[-\frac{k^2c_\nu^2}{2}\langle Q_\nu^2\rangle_T\right]
 =\exp\!\left[-\frac{k^2}{2}\sum_\nu c_\nu^2
 \langle Q_\nu^2\rangle_T\right].
\end{equation*}
После изотропного усреднения это даёт ранее найденную формулу
$\mathcal A_T=e^{-k^2\langle u^2\rangle_T/6}$.

\clearpage
\section{Фактор Лэмба--Мёссбауэра в модели Дебая}
Переходим от суммы по модам к плотности состояний
$D_D(\omega)=9N\omega^2/\omega_D^3$, $0<\omega<\omega_D$:
\begin{align*}
 \langle u^2\rangle_T
 &=\frac\hbar{NM}\int_0^{\omega_D}\frac{D_D(\omega)}\omega
 \left(\frac12+\frac1{e^{\hbar\omega/T}-1}\right)d\omega\\
 &=\frac{9\hbar}{M\omega_D^3}
 \left[\frac{\omega_D^2}{4}
 +\left(\frac T\hbar\right)^2
 \int_0^{\Theta_D/T}\frac{z\,dz}{e^z-1}\right].
\end{align*}
Подставим это выражение в $f_{\mathrm{LM}}=e^{-k^2\langle u^2\rangle_T/3}$.
С учётом $R=\hbar^2k^2/(2M)$ и $\Theta_D=\hbar\omega_D$ имеем
\begin{equation*}
 \boxed{f_{\mathrm{LM}}(T)=\exp\!\left\{
 -\frac{6R}{\Theta_D}\left[\frac14+
 \left(\frac{T}{\Theta_D}\right)^2
 \int_0^{\Theta_D/T}\frac{z\,dz}{e^z-1}\right]\right\}.}
\end{equation*}

\subsection{Низкие температуры}
При $T\ll\Theta_D$ верхний предел интеграла можно заменить бесконечностью:
\begin{equation*}
 \int_0^\infty\frac{z\,dz}{e^z-1}=\Gamma(2)\zeta(2)=\frac{\pi^2}{6}.
\end{equation*}
Отсюда
\begin{equation*}
 f_{\mathrm{LM}}(T)\simeq
 \exp\!\left[-\frac{3R}{2\Theta_D}
 \left(1+\frac{2\pi^2}{3}\frac{T^2}{\Theta_D^2}\right)\right],
 \qquad f_{\mathrm{LM}}(0)=e^{-3R/(2\Theta_D)}.
\end{equation*}

\subsection{Высокие температуры}
При $T\gg\Theta_D$ используем
$\langle u^2\rangle_T\simeq9T/(M\omega_D^2)$:
\begin{equation*}
 f_{\mathrm{LM}}(T)\simeq
 \exp\!\left[-\frac{3k^2T}{M\omega_D^2}\right]
 =\exp\!\left[-\frac{6RT}{\Theta_D^2}\right].
\end{equation*}
Повышение температуры уменьшает бесфононную долю.
Поэтому для наблюдения эффекта благоприятны малый $R$, высокая
характерная частота решётки и низкая температура.

\subsection{Численная оценка для железа}
Положим $\Theta_D/k_B=430\,\text{К}$,
$\Theta_D\simeq0.0371\,\text{эВ}$ и $R\simeq0.00195\,\text{эВ}$.
Полная дебаевская формула даёт
\begin{equation*}
 f_{\mathrm{LM}}(0)\simeq0.924,\qquad
 f_{\mathrm{LM}}(77\,\text{К})\simeq0.909,\qquad
 f_{\mathrm{LM}}(300\,\text{К})\simeq0.792.
\end{equation*}
Здесь аргументы в кельвинах обозначают физическую температуру $T_{\mathrm K}$.
Оценки относятся к одноатомной изотропной модели Дебая с указанным
параметром; реальная вероятность зависит от колебательного спектра
и окружения ядра.

\end{document}
